\subsection{Characterization of integral elements by means of places}

PROPOSITION 5. Let K be a field, S a subring of K, h a homomorphism from S to a field and p the kernel of h. For an element x of K to be integral over the local ring \(S_{p}\) , it is necessary and sufficient that every place of K extending h be finite at x.

If \(f\) is a place of \(K\) extending \(h, f\) is finite on \(S_p\) and zero on \(pS_p\) and hence the ring of the place \(f\) dominates \(S_p\) . Conversely, if \(V\) is a valuation ring of \(K\) which dominates \(S_p\) , \(V\) is the ring of a place \(f\) whose restriction to \(S\) is a homomorphism with the same kernel as \(h\) ; replacing \(f\) by an equivalent place, it is seen that \(V\) is the ring of a place of \(K\) which extends \(h\) . To say that every place of \(K\) extending \(h\) is finite at \(x\) is equivalent to saying that \(x\) belongs to all the valuation rings of \(K\) which dominate \(S_p\) . The proposition then follows from Theorem 3 of § 1, no. 3.

PROPOSITION 6. Let K be a field and S a subring of K. For an element \(x \in \mathbf{K}\) to be integral over S, it is necessary and sufficient that every place of K which is finite on S be finite at \(x\) .

This is also a consequence of Theorem 3 of §1, no. 3.

